\MajorQuestionLesson{一般動詞の現在形と肯定文}

\TypeQuestion{英文を書き写し、日本語に訳しなさい。}{generalverb_affirmative_copy_translate}
\EngCopyTranslateSplit{I play tennis.}{私はテニスをします。}
\EngCopyTranslateSplit{You like music.}{あなたは音楽が好きです。}
\EngCopyTranslateSplit{We use this room after school.}{私たちは放課後にこの部屋を使います。}
\EngCopyTranslateSplit{They study English every day.}{彼らは毎日英語を勉強します。}
\EngCopyTranslateSplit{I am busy.}{私は忙しいです。}
\EngCopyTranslateSplit{You are my friend.}{あなたは私の友達です。}

\LessonPageBreak

\TypeQuestion{英文を書き写したあと、主語に下線を引き、動詞を丸で囲みなさい。}{generalverb_affirmative_copy_structure}
\SetRowSpace{3mm}
\EngWriteQuestion{\EngSentence{I like music.}}{\underline{I} \EngCircle{like} music.}
\EngWriteQuestion{\EngSentence{They play soccer after school.}}{\underline{They} \EngCircle{play} soccer after school.}
\EngWriteQuestion{\EngSentence{I am a student.}}{\underline{I} \EngCircle{am} a student.}
\EngWriteQuestion{\EngSentence{You are kind.}}{\underline{You} \EngCircle{are} kind.}
\EngWriteQuestion{\EngSentence{We use this room.}}{\underline{We} \EngCircle{use} this room.}
\EngWriteQuestion{\EngSentence{They are busy.}}{\underline{They} \EngCircle{are} busy.}
\EngWriteQuestion{\EngSentence{You have a pen.}}{\underline{You} \EngCircle{have} a pen.}
\EngWriteQuestion{\EngSentence{We read books at school.}}{\underline{We} \EngCircle{read} books at school.}
\EngWriteQuestion{\EngSentence{This is my book.}}{\underline{This} \EngCircle{is} my book.}
\EngWriteQuestion{\EngSentence{I get up at seven.}}{\underline{I} \EngCircle{get} up at seven.}

\LessonPageBreak

\TypeQuestion{日本語に合うように、英文を必要なところだけ直して書きなさい。\par{\small ※短縮形の有無は、もとの英文に合わせること。}}{generalverb_affirmative_error_correct}
\SetRowSpace{3mm}
\EngWriteQuestion{\EngSentence{I am play tennis.}\quad〔私はテニスをします。〕}{I play tennis.}
\EngWriteQuestion{\EngSentence{We are study English.}\quad〔私たちは英語を勉強します。〕}{We study English.}
\EngWriteQuestion{\EngSentence{i like music.}\quad〔私は音楽が好きです。〕}{I like music.}
\EngWriteQuestion{\EngSentence{They play soccer}\quad〔彼らはサッカーをします。〕}{They play soccer.}
\EngWriteQuestion{\EngSentence{I a pen have.}\quad〔私はペンを持っています。〕}{I have a pen.}
\EngWriteQuestion{\EngSentence{You am busy.}\quad〔あなたは忙しいです。〕}{You are busy.}
\EngWriteQuestion{\EngSentence{They is my friends.}\quad〔彼らは私の友達です。〕}{They are my friends.}
\EngWriteQuestion{\EngSentence{You this book read.}\quad〔あなたはこの本を読みます。〕}{You read this book.}
\EngWriteQuestion{\EngSentence{We like music}\quad〔私たちは音楽が好きです。〕}{We like music.}
\EngWriteQuestion{\EngSentence{This are my pen.}\quad〔これは私のペンです。〕}{This is my pen.}

\LessonPageBreak

\TypeQuestion{日本語に合うように、〔　〕内のすべての語句を並べ替えて、正しい英文を作りなさい。}{generalverb_affirmative_reorder}
\SetRowSpace{2mm}
\EngWriteQuestion{私は7時に起きます。\par\EngReorder{at seven / get up / i}}{I get up at seven.}
\EngWriteQuestion{あなたは音楽が好きです。\par\EngReorder{music / you / like}}{You like music.}
\EngWriteQuestion{私は生徒です。\par\EngReorder{a student / am / i}}{I am a student.}
\EngWriteQuestion{私たちは学校で昼食を食べます。\par\EngReorder{at school / have lunch / we}}{We have lunch at school.}
\EngWriteQuestion{彼らは毎日英語を勉強します。\par\EngReorder{every day / English / they / study}}{They study English every day.}
\EngWriteQuestion{あなたは忙しいです。\par\EngReorder{busy / are / you}}{You are busy.}
\EngWriteQuestion{これは私の本です。\par\EngReorder{my book / this / is}}{This is my book.}
\EngWriteQuestion{私はペンを持っています。\par\EngReorder{have / i / a pen}}{I have a pen.}

\WordBank{一般動詞の現在形と肯定文}{
  \WB{I}{代名詞}{私は}{必ず大文字で書く}
  \WB{you}{代名詞}{あなた（たち）は}{}
  \WB{we}{代名詞}{私たちは}{}
  \WB{they}{代名詞}{彼ら・彼女ら・それらは}{}
  \WB{my}{所有格}{私の}{}
  \WB{this}{代名詞等}{これ；この}{}
  \WB{a}{冠詞}{1つの／1人の}{訳さない場合もある}
  \WB{am / are / is}{be動詞}{～です／～にいます}{主語に合わせて使い分ける}
  \WB{play}{動詞}{（スポーツを）する}{}
  \WB{study}{動詞}{勉強する}{}
  \WB{like}{動詞}{～が好きである}{}
  \WB{use}{動詞}{使う}{}
  \WB{have}{動詞}{持っている；食べる}{}
  \WB{read}{動詞}{読む}{}
  \WB{get}{動詞}{得る}{\EngText{get up}で「起きる」}
  \WB{music}{名詞}{音楽}{}
  \WB{English}{名詞}{英語}{最初の文字は大文字}
  \WB{tennis}{名詞}{テニス}{}
  \WB{soccer}{名詞}{サッカー}{}
  \WB{book}{名詞}{本}{複数形 \EngText{books}}
  \WB{pen}{名詞}{ペン}{}
  \WB{student}{名詞}{生徒}{複数形 \EngText{students}}
  \WB{friend}{名詞}{友達}{複数形 \EngText{friends}}
  \WB{room}{名詞}{部屋}{}
  \WB{school}{名詞}{学校}{}
  \WB{lunch}{名詞}{昼食}{}
  \WB{day}{名詞}{日；1日}{}
  \WB{busy}{形容詞}{忙しい}{}
  \WB{kind}{形容詞}{親切な}{}
  \WB{every}{形容詞}{毎～}{}
  \WB{seven}{数詞}{7（七）}{}
  \WB{at}{前置詞}{～で／～時に}{}
  \WB{after}{前置詞}{～のあとに}{}
  \WB{up}{副詞}{上に}{}
  \WB{get up}{熟語}{起きる}{}
  \WB{have lunch}{熟語}{昼食を食べる}{}
  \WB{at school}{語句}{学校で}{}
  \WB{at seven}{語句}{7時に}{}
  \WB{after school}{熟語}{放課後に}{}
  \WB{every day}{熟語}{毎日}{}
}

\PhraseBank{一般動詞の現在形と肯定文}{
  \PB{I play tennis.}{私はテニスをします。}
  \PB{You like music.}{あなたは音楽が好きです。}
  \PB{We study English every day.}{私たちは毎日英語を勉強します。}
  \PB{They use this room.}{彼らはこの部屋を使います。}
}
